% ===== §12 Dissipative Structure and the Thermodynamics of Generativity =====
\section{Dissipative Structure and the Thermodynamics of Generativity}
\label{sec:dissipation}

\noindent
The three questions have their three answers, and a structural account of power under intimacy is, in its essentials, complete. What remains is to set that account within the widest frame the series commands---to show that the findings reached by the analysis of code, prediction, value, power, and history are not local peculiarities of the human bond but instances of something that holds of generative systems as such. This is the office of the three cosmological sections. The first establishes the thermodynamic frame: that a generative system is a system held far from equilibrium, that estrangement is the necessary signature of its being so held, and that the whole etiology of the diagnostic chapters is, seen from here, the microphysics of a single macroscopic necessity.

\subsection{Far from equilibrium}
\label{sub:dissipation-far}

The decisive result of non-equilibrium thermodynamics is that order, in certain systems, arises not despite but \emph{because} of their being driven far from equilibrium. A dissipative structure---a convection cell, a chemical clock, a living organism---maintains its order only by a continual throughput of energy and a continual export of entropy; it exists as an ordered thing exactly insofar as it is held away from the equilibrium toward which an isolated system would relax.\footnote{For dissipative structures and the thesis that order can emerge and be sustained only in systems maintained far from thermodynamic equilibrium by a continual flux, see \citet{prigogine1984}.} Equilibrium, for such a system, is not rest but death: the state of maximum entropy, of no gradient, no flux, no further change---the state a living thing reaches when it stops being a living thing.

The generative relation is a dissipative structure in this exact sense. Its order---the ongoing generation of relational value---is sustained only by a continual throughput: of difference, of surprise, of the prediction-error that is the nutrition of coupling (\xref{sec:prediction}), of the contradiction that the dialectic works into new form (\xref{sec:histmat}). Cut off the throughput, resolve all difference, settle every contradiction, and the relation does not achieve a blissful stability; it reaches relational equilibrium, which is the heat-death of the bond. Here the political vocabulary of the negative summit and the thermodynamic vocabulary converge exactly: \emph{static equality is the political name of thermodynamic equilibrium}. The fixed, settled, frictionless arrangement toward which the optimist aspires is the state of no gradient and no flux---and a relation in that state generates nothing, because generation \emph{is} the throughput, and the throughput has stopped.

\begin{claim}[The relation as dissipative structure]
A generative relation maintains its order---the ongoing generation of value---only as a system held far from equilibrium, by a continual throughput of difference, surprise, and contradiction. Equilibrium is not its rest but its death: the frictionless concordance toward which the optimist aspires is the state of no gradient and no flux, in which nothing is generated because generation \emph{is} the throughput. Static equality is the political name of thermodynamic equilibrium; and the perpetual non-settling that the negative summit demanded is the thermodynamic condition of remaining alive.
\end{claim}

\subsection{Estrangement as the signature of disequilibrium}
\label{sub:dissipation-signature}

This frame transfigures the paper's first claim and gives it its deepest statement. If a generative relation is a system held far from equilibrium, then it cannot be smooth. A far-from-equilibrium system fluctuates; it passes through phases; it does not hold a single steady state but cycles, and the cycling is the very mode of its persistence. \emph{Surechigai}---the passing-without-meeting---is, at this level, nothing other than the \emph{signature of disequilibrium}: the visible mark that the system is still being driven, still fluctuating, still alive. A relation that never passed, that held perfect and unbroken concordance, would be displaying the one thing a living dissipative structure cannot display: a steady state at equilibrium. The absence of all estrangement would not be the perfection of the bond but the diagnostic sign of its death.

This is why the symptomatology insisted (\xref{sub:symptom-distinction}) that not every symptom is a pathology, and why the cardinal distinction was between reparable misalignment and structural shadow. We can now say precisely what the structural shadow \emph{is}: it is the estrangement that is not a fault in the system but the necessary fluctuation of a system held from equilibrium---the passing that must occur because the bond is alive and being driven. To attack such estrangement as a malfunction is to attack the disequilibrium itself, which is to push the system toward the equilibrium that is its death. The most destructive misreading in the whole etiology is the demand that a living bond stop fluctuating---that it achieve, once and for all, the steady concordance whose name, correctly given, is heat-death.

\subsection{The etiology as microphysics}
\label{sub:dissipation-microphysics}

The thermodynamic frame also retrospectively unifies the diagnostic chapters. The four causes traced through code, prediction, value, and power looked, in their own chapters, like four distinct mechanisms of breakdown. Seen from the thermodynamic level, they are the \emph{microphysics of the dissipative phase}---the local, mechanistic detail of how a far-from-equilibrium system fluctuates. The mismatch of grammars, the unattributable error and predictive withdrawal, the arrest of phase, the seepage of peripheral power: these are the channels through which the necessary fluctuation actually propagates, the specific machinery by which a generative system, being driven, passes through its disequilibrium phases. The etiology does not compete with the thermodynamic account; it is its fine grain. And this dissolves what might have seemed a tension in the paper---between an etiology that treats estrangement as caused, layer by layer, and a functional claim that treats it as necessary. The two are the same fact at two scales: the necessity is macroscopic (a driven system must fluctuate), and the causation is microscopic (this is the machinery through which it fluctuates here). Cause and necessity are not rivals but levels.

\begin{claim}[Etiology and function are levels, not rivals]
The layered causation of estrangement (the etiology) and its structural necessity (the functional claim) are not in tension but are one fact at two scales. Macroscopically, a generative system held far from equilibrium \emph{must} fluctuate---estrangement is necessary. Microscopically, the four causes are the machinery through which the fluctuation propagates in a human bond---estrangement is caused, layer by layer. The etiology is the microphysics of the dissipative phase; the functional necessity is its macrophysics. To have an account of how estrangement is caused is not to have refuted the claim that it is necessary, any more than knowing the mechanism of a fluctuation refutes the law that the driven system must fluctuate.
\end{claim}

\subsection{The placement of necessity}
\label{sub:dissipation-placement}

A word, finally, on where in the paper's scheme this necessity belongs---a matter left open in the catalogue of causes. The thermodynamic necessity of estrangement was listed, provisionally, among the nineteen causes, as the structural-level entry that explained not any particular passing but why passing must occur at all. It is now clear that it does not sit comfortably there, for it is not a cause in the sense the other entries are. The mismatch of grammars, the unattributable error, the seepage of power---each is a mechanism that produces a particular estrangement, and might in a given case be absent. The thermodynamic necessity is of a different kind: it is not a mechanism that might or might not operate but the macroscopic law that \emph{some} mechanism must, because a living relation is a driven system. It is properly not the nineteenth cause but the \emph{functional ground} of there being causes at all---the reason the etiology has work to do. The paper therefore lifts it out of the catalogue of mechanisms and lodges it here, in the thermodynamic frame, where it belongs: not as one more cause of estrangement, but as the law that makes estrangement, in some form, inevitable for anything alive enough to generate.

\noindent
With the thermodynamic frame established---disequilibrium as the condition of generative life, estrangement as its signature, the etiology as its microphysics---the cosmological movement can take its next step, from the physics of the driven system to the grammar of its return. For a system that fluctuates does not merely depart; it comes back, and the form of its coming-back is the subject of the oldest cosmology the series draws upon.
